\section{The full matrix range behind the graph problem}\label{ah:sec:araki-horn}
All matrices are real. Spectra and vectors are listed in decreasing
order. For real vectors $x,y\in\mathbb R^n$, write $x\prec y$ when
$\sum_{i=1}^k x_i\le\sum_{i=1}^k y_i$ for $k<n$, with equality
for $k=n$. For positive vectors, $q\prec_{\log}p$ means
$\log q\prec\log p$.

\begin{theorem}[Relative spectral closure]\label{ah:thm:range}
Fix $n\ge1$, $0<\alpha<1$, and a decreasing list
$r_1\ge\cdots\ge r_n>0$. Let $\mathcal R_{\alpha,r}$ consist of
the ordered spectra of
\[
 Q_\alpha(A,B)=A^{-\alpha/2}B^\alpha A^{-\alpha/2},
 \qquad A,B\succ0,
 \qquad \spec(A^{-1/2}BA^{-1/2})=r.
\]
Then
\[
 \overline{\mathcal R_{\alpha,r}}
 =\{q_1\ge\cdots\ge q_n>0:\log q\prec\alpha\log r\}.
\]
The same assertion holds if $A$ is required to be diagonal with entries
$\delta^{n-i}$ and $\delta\downarrow0$ is allowed in the closure.
If $r_i\ge1$ for every $i$, the witnesses satisfy $B\succeq A$ and
$\rank(B-A)=\#\{i:r_i>1\}$, exactly.
\end{theorem}

The last rank statement follows because a congruence of a matrix with
eigenvalues $r_i-1$ has the same inertia. Thus the theorem preserves
the complete relative input spectrum, not merely its determinant.
In the rank-one case $r=(1+\tau,1,\ldots,1)$, majorization is
equivalent to $q_i\ge1$ and $\prod_iq_i=(1+\tau)^\alpha$.

\begin{lemma}[General reverse-Schur fractional limit]\label{ah:lem:pivots}
Let $S\succ0$ be fixed. Define
\[
 a_i=\delta^{n-i},\quad A_\delta=\diag(a_i),\quad
 B_\delta=A_\delta^{1/2}S A_\delta^{1/2}.
\]
Let $d_i>0$ be the reverse Schur pivots of $S$:
\[
 d_i=\frac{\det S[i:n,i:n]}{\det S[i+1:n,i+1:n]},
 \qquad\det S[n+1:n,n+1:n]=1.
\]
Then
\[
 A_\delta^{-\alpha/2}B_\delta^\alpha A_\delta^{-\alpha/2}
 \longrightarrow\diag(d_1^\alpha,\ldots,d_n^\alpha)
 \quad\text{in operator norm}.
\]
Convergence is uniform when $S$ ranges over a compact subset of the
SPD cone.
\end{lemma}

\begin{proof}
Choose constants $0<m\le M$ with $mI\preceq S\preceq MI$.
Write $\nu_1\le\cdots\le\nu_n$ for the eigenvalues of $B_\delta$.
The min--max principle gives
\[
 ma_i\le\nu_i\le Ma_i.
\]
In particular distinct scales imply $\nu_i/\nu_{i+1}\to0$.
For $1\le k\le n$, the elementary symmetric polynomial satisfies
\[
 e_k(B_\delta)
 =\sum_{|I|=k}\det S[I,I]\prod_{i\in I}a_i
 \sim\det S[n-k+1:n,n-k+1:n]\prod_{i=n-k+1}^n a_i.
\]
Every other product has an extra positive power of $\delta$.
Likewise, separation of the $\nu_i$ implies
$e_k(B_\delta)\sim\prod_{i=n-k+1}^n\nu_i$.
Taking successive quotients shows
\[
 \nu_i/a_i\longrightarrow d_i.
\]

Let $v_k$ be orthonormal eigenvectors associated with $\nu_k$,
and write $v_{ik}$ for their coordinates. Positivity of $B_\delta$
and $B_\delta^{-1}$ gives two estimates:
\[
 |v_{ik}|^2\le\frac{(B_\delta)_{ii}}{\nu_k}
 \le C\frac{a_i}{a_k},\qquad
 |v_{ik}|^2\le\nu_k(B_\delta^{-1})_{ii}
 \le C\frac{a_k}{a_i}.
\]
Thus $|v_{ik}|^2\le C\min(a_i/a_k,a_k/a_i)$, and
$|v_{ii}|^2\to1$. Define the nonnegative summands
\[
 t_{ik}=a_i^{-\alpha}\nu_k^\alpha|v_{ik}|^2.
\]
For $k>i$, $t_{ik}=O((a_i/a_k)^{1-\alpha})$; for $k<i$,
$t_{ik}=O((a_k/a_i)^{1+\alpha})$. Therefore
$t_{ii}\to d_i^\alpha$ and $t_{ik}\to0$ whenever $i\ne k$.
These are the diagonal summands of the relative fractional matrix.
For its off-diagonal entries,
\[
 |(Q_\alpha)_{ij}|\le\sum_k\sqrt{t_{ik}t_{jk}}\to0
 \qquad(i\ne j).
\]
For each $k$ at least one factor vanishes and all factors stay bounded.
This proves entrywise, hence norm, convergence in fixed dimension.
On compact SPD sets, $m,M,C$ and the nonzero determinant lower bounds
can all be chosen uniformly, proving the last assertion.
\end{proof}

The successive-Schur asymptotics and eigenvector scale bounds have prior
sources: Usevich--Barthelm\'e~\cite[Theorem 4.2 in the arXiv v2 PDF]{UB} and
Gissler--Auger--Hansen~\cite[Theorem 1]{GAH}. The proof above is included
to check that the power normalization permits precisely $0<\alpha<1$;
the scaling mechanism is not presented as new.

\begin{proof}[Proof of Theorem~\ref{ah:thm:range}]
For $X,Y\succ0$, Araki log-majorization~\cite{Araki} states
\[
 \lambda(X^{\alpha/2}Y^\alpha X^{\alpha/2})
 \prec_{\log}
 \lambda(X^{1/2}YX^{1/2})^\alpha,
 \qquad0<\alpha<1.
\]
Apply this with $X=A^{-1}$ and $Y=B$. This proves necessity.
All admissible $q_i$ lie in the compact interval
$[r_n^\alpha,r_1^\alpha]$, so the necessary set is closed.

Conversely suppose $\log q\prec\alpha\log r$, and put
$d_i=q_i^{1/\alpha}$. The real triangular version of the
Weyl--Horn theorem~\cite{Horn,LiMathias} produces a real upper
triangular matrix $R$ with positive diagonal $\sqrt{d_i}$ and
singular values $\sqrt{r_i}$. Set $S=RR^T$.
It has eigenvalues $r_i$, and its reverse Schur pivots are exactly
$d_i$. Indeed the trailing principal block is the corresponding
trailing triangular factor times its transpose, so its determinant
is the product of those diagonal squares.
Lemma~\ref{ah:lem:pivots} now provides the required witnesses, and
$A_\delta^{-1/2}B_\delta A_\delta^{-1/2}=S$ at every $\delta$.
The rank assertion follows from Sylvester's law of inertia.
\end{proof}

Huang--Tam~\cite[Theorem 3.3]{HT} explicitly characterize squared
Cholesky diagonals on a unitary orbit by multiplicative majorization,
using Kostant convexity~\cite{Kostant}. We use the real Weyl--Horn
construction above to remove any possible real-versus-complex issue.

\subsection{Sharp envelopes and counts}
\begin{corollary}[Every individual spectral envelope is sharp]
\label{ah:cor:envelope}
Under the hypotheses of Theorem~\ref{ah:thm:range}, for $1\le k\le n$,
\[
 \sup\lambda_k(Q_\alpha)
 =\left(\prod_{i=1}^k r_i\right)^{\alpha/k},
 \qquad
 \inf\lambda_{n-k+1}(Q_\alpha)
 =\left(\prod_{i=n-k+1}^n r_i\right)^{\alpha/k}.
\]
The extrema are suprema or infima over actual SPD pairs; attainment
at finite condition number is not asserted.
\end{corollary}
\begin{proof}
If $x=\log\lambda(Q_\alpha)$ and $y=\alpha\log r$, then
$kx_k\le\sum_{i=1}^kx_i\le\sum_{i=1}^ky_i$.
To saturate the first bound in the closure, replace the first $k$
entries of $y$ by their mean and keep all other entries unchanged.
The resulting ordered vector is majorized by $y$.
Average the bottom $k$ entries for the lower bound.
\end{proof}

\begin{corollary}[Exact strict counts for a fixed input spectrum]
\label{ah:cor:count}
For $s>0$,
\[
 \sup\#\{j:\lambda_j(Q_\alpha)>s\}
 =\max\left(\{0\}\cup
 \left\{k\in\{1,\ldots,n\}:
 \alpha\sum_{i=1}^k\log r_i>k\log s\right\}\right).
\]
Likewise
\[
 \sup\#\{j:\lambda_j(Q_\alpha)<s\}
 =\max\left(\{0\}\cup
 \left\{k\in\{1,\ldots,n\}:
 \alpha\sum_{i=n-k+1}^n\log r_i<k\log s\right\}\right).
\]
\end{corollary}
\begin{proof}
The envelope bounds give necessity. If the relevant inequality is
strict, the averaged target spectrum in the preceding proof has
the required $k$ eigenvalues strictly on the desired side. A sufficiently
close finite witness preserves the strict inequalities.
\end{proof}

For a rank-one addition with leverage $\tau$, choose $s=1+\varepsilon$.
For $k\ge1$ the condition is simply
$k\log(1+\varepsilon)<\alpha\log(1+\tau)$.
This recovers the strict ceiling-minus-one formula in Corollary~\ref{mat:cor:count}. For a batch, the first few eigenvalue envelopes can be
strictly stronger than a bound that uses only the determinant.

\subsection{The effective-resistance Gram matrix of a batch}
Let $L$ be a connected graph Laplacian on $N$ vertices, and let
\[
 L'=L+UU^T,\qquad
 U=(\sqrt{w_1}b_1,\ldots,\sqrt{w_m}b_m),\quad w_j>0,
\]
where the $b_j$ are edge-incidence vectors.
All matrices in this section act on $\mathbf1^\perp$, of dimension
$d=N-1$. Let $\gamma_1\ge\cdots\ge\gamma_r>0$ be the positive
eigenvalues of the effective-resistance Gram matrix
\[
 G=U^TL^\dagger U,
\]
and pad $\gamma_i=0$ to length $d$. Then the original relative
input eigenvalues are $r_i=1+\gamma_i$.

\begin{corollary}[A complete-Gram spectral envelope]\label{ah:cor:graph}
For
\[
 C=(L^\dagger)^{\alpha/2}(L'^\alpha-L^\alpha)
      (L^\dagger)^{\alpha/2}
\]
on $\mathbf1^\perp$, one has for $1\le k\le d$
\[
 0\le\lambda_k(C)
 \le\exp\left(\frac\alpha k
      \sum_{i=1}^{\min(k,r)}\log(1+\gamma_i)\right)-1.
\]
Consequently a correction $F$ of rank at most $k$ exists with
\[
 0\preceq (L'^\alpha-L^\alpha)-F\preceq\varepsilon L^\alpha
\]
as soon as $k<d$ and
\[
 \alpha\sum_{i=1}^{\min(k+1,r)}\log(1+\gamma_i)
 \le(k+1)\log(1+\varepsilon).
\]
The exact universal matrix worst-case count for this fixed list
$\gamma$ is Corollary~\ref{ah:cor:count} with $s=1+\varepsilon$.
For graph additions the bound always holds; for one edited edge its
unrestricted sharpness on weighted unit-mass paths follows from Theorem~\ref{w:thm:path}.
\end{corollary}
\begin{proof}
The nonunit eigenvalues of $L^{-1/2}L'L^{-1/2}$ on the common support
are $1+\gamma_i$. Apply Corollary~\ref{ah:cor:envelope}.
Truncate $C$ at its largest $k$ eigenvalues, obtaining $C_k$,
and set $F=L^{\alpha/2}C_kL^{\alpha/2}$.
Then $0\preceq C-C_k\preceq\lambda_{k+1}(C)I$, which proves the
energy inequality after congruence.
\end{proof}

This is an existence and obstruction statement for low-rank
preconditioner updates. It does not give a fast method to compute the
spectral truncation. The small Gram matrix supplies a checkable rank
certificate, but computing it requires effective-resistance solves.
No realization of every higher-rank fixed-$\gamma$ extremum by a
prescribed graph class is claimed here.

